\documentclass[11pt]{article}
\usepackage[T1]{fontenc}
\usepackage{lmodern}
\usepackage[margin=1in]{geometry}
\usepackage{amsmath,amssymb,amsthm,mathtools}
\usepackage{microtype,booktabs,array}
\usepackage{verbatim}
\usepackage[colorlinks=true,linkcolor=blue,citecolor=blue,urlcolor=blue]{hyperref}
\hypersetup{pdftitle={An explicit exponent above 65/32 for ordinary sensitivity versus block sensitivity},pdfauthor={ChatGPT},pdfsubject={Indexed tournament recurrences and an exact integer separation certificate}}
\newtheorem{theorem}{Theorem}[section]
\newtheorem{lemma}[theorem]{Lemma}
\newtheorem{proposition}[theorem]{Proposition}
\newtheorem{corollary}[theorem]{Corollary}
\theoremstyle{definition}
\newtheorem{definition}[theorem]{Definition}
\theoremstyle{remark}
\newtheorem{remark}[theorem]{Remark}
\DeclareMathOperator{\bs}{bs}
\DeclareMathOperator{\dist}{dist}
\newcommand{\sens}{s}
\newcommand{\bits}{\{0,1\}}
\newcommand{\zero}{\mathbf{0}}
\newcommand{\one}{\mathbf{1}}
\newcommand{\NN}{\mathbb{N}}
\newcommand{\RR}{\mathbb{R}}
\newcommand{\floor}[1]{\left\lfloor #1\right\rfloor}
\newcommand{\ceil}[1]{\left\lceil #1\right\rceil}
\setlength{\emergencystretch}{2em}

\title{An explicit exponent above $65/32$\\
for ordinary sensitivity versus block sensitivity}
\author{Research companion prepared with ChatGPT\\for the ProveIt repository}
\date{7 October 2026}

\begin{document}
\maketitle

\begin{abstract}
We give an explicitly quantified refinement of the nested-tournament
construction in the repository manuscript \emph{A superquadratic separation
between sensitivity and block sensitivity}.  A fixed total Boolean function
has sensitivity at most $S_*=9\cdot2^{318}$ and block sensitivity at zero at
least $B_*=100\cdot2^{646}$.  Its iterated disjoint compositions consequently
satisfy
\[
 \bs(F_m,\zero)\ge \sens(F_m)^{\alpha_*},\qquad
 \alpha_*=
 \frac{\log100+646\log2}{\log9+318\log2}
 =2.03208272439\ldots>\frac{65}{32}.
\]
Moreover, $\bs(F_m,\zero)/\sens(F_m)^{65/32}$ tends to infinity.
The refinement allows different tournament sizes at different depths,
retains the sensitivity of each nested predicate separately, and sharpens
the count of clauses that are one gate away from accepting.  We prove all
combinatorial and composition statements used here.  The remaining finite
calculation consists of twenty applications of explicit integer
recurrences; a complete dependency-free verifier is included.  The
construction is an existence theorem with a finite search procedure for
its labels, not a practical truth-table construction.  No optimality of
the exponent, and no new Lean formalization, is claimed.
\end{abstract}

\tableofcontents

\section{Main result, attribution, and scope}
\label{sec:main}

Ordinary sensitivity counts input bits that individually change a Boolean
function's value.  Block sensitivity permits disjoint nonempty groups of
bits to change together.  Huang's theorem settled the polynomial relation
between these measures \cite{Huang2019}; determining the best possible
exponent remains a finer quantitative problem.  It is essential here to
use ordinary sensitivity.  A separation involving the smaller spectral
sensitivity does not by itself give the result below.

The starting point is OpenAI's September 25, 2026 repository manuscript
\cite{OpenAI2026}.  It adapts the tournament architecture of Meiburg
\cite{Meiburg2026} to nested predicates and proves that the ratio
$\bs(f)/\sens(f)^2$ is unbounded.  We inspected the four source sections
giving the construction, recurrences, numerical bounds, and composition
argument at the pinned repository revision recorded in the references.
The repository also supplies a Lean proof entry and a comparator
configuration for its quantitative separation theorem.  That original
formalization was not rebuilt for this companion; the new estimates below
are proved in prose and checked by exact arithmetic.

\begin{theorem}[Explicit ordinary-sensitivity exponent]
\label{thm:main}
Put
\[
 S_*=9\cdot2^{318},\qquad B_*=100\cdot2^{646},\qquad
 \alpha_*=\frac{\log B_*}{\log S_*}.
\]
There is a positive integer $n_*$ and, for every integer $m\ge1$, a
nonconstant total Boolean function
$F_m:\bits^{n_*^m}\to\bits$ such that
\begin{gather}
 F_m(\zero)=0,\qquad
 \sens(F_m)\le S_*^m,\qquad
 \bs(F_m,\zero)\ge B_*^m,
 \label{eq:main-bounds}\\
 \bs(F_m,\zero)\ge \sens(F_m)^{\alpha_*},\qquad
 \alpha_*>\frac{65}{32}.
 \label{eq:main-exponent}
\end{gather}
In particular,
\[
 \frac{\bs(F_m,\zero)}{\sens(F_m)^{65/32}}
 \longrightarrow\infty.
\]
The number $n_*$ is specified by Equation~\eqref{eq:seed-dimension}.
\end{theorem}

The quantitative comparison with the source deserves precision.  Its
printed fixed-power example takes depth $d=9$ and $M=9^{10}$, giving an
exponent approximately $2.00025240714$.  The same published inequalities,
without any new lemma, already give approximately $2.00647082823$ when
one instead takes $d=20$ and $M=9^{21}$.  Our result therefore should not be
described merely as an improvement over the source's first displayed
seed.  We change the parameter schedule and its analysis, obtaining the
larger certified exponent in Theorem~\ref{thm:main}.

\begin{center}
\begin{tabular}{@{}p{0.55\linewidth}l@{}}
\toprule
Certificate being compared & Certified exponent\\
\midrule
Source's displayed seed, $d=9$, $M=9^{10}$
  & $2.00025240714\ldots$\\
Source's same estimate evaluated at $d=20$, $M=9^{21}$
  & $2.00647082823\ldots$\\
Theorem~\ref{thm:main}, using the simple bounds $S_*,B_*$
  & $2.03208272439\ldots$\\
\bottomrule
\end{tabular}
\end{center}

These numerical values describe particular sufficient estimates.  We do
not prove that the source's parameter family has no better choice, nor
that our schedule is optimal within its more general family.  Likewise,
the finite label search implicit in the proof is far too large for a
practical construction.  Its purpose is to establish an unambiguous
finite Boolean function with controlled complexity measures.

\section{Definitions and bookkeeping}

For a positive integer $n$, write $[n]=\{1,\ldots,n\}$.  If
$x\in\bits^n$ and $B\subseteq[n]$, let $x^B$ be obtained from $x$ by
flipping precisely the coordinates in $B$.  Write $x^a=x^{\{a\}}$.
For a total function $f:\bits^n\to\bits$, define
\begin{align*}
 \sens(f,x)&=|\{a\in[n]:f(x^a)\ne f(x)\}|,\\
 \sens(f)&=\max_{x\in\bits^n}\sens(f,x).
\end{align*}
The block sensitivity $\bs(f,x)$ is the largest cardinality of a family
of pairwise disjoint nonempty sets $B_1,\ldots,B_b\subseteq[n]$ for
which $f(x^{B_j})\ne f(x)$ for every $j$; put
$\bs(f)=\max_x\bs(f,x)$.

We will construct a family of predicates $P_1,\ldots,P_H$ on the same
input space satisfying
\[
                    P_1(x)\ge P_2(x)\ge\cdots\ge P_H(x)
                    \quad\text{for all }x.
\]
Nesting refers to the predicate index and does not assert monotonicity
in the input bits.  Let
\begin{align*}
 S_0(q)&=\max_{x:P_q(x)=0}|\{a:P_q(x^a)=1\}|,\\
 S_1(q)&=\max_{x:P_q(x)=1}|\{a:P_q(x^a)=0\}|.
\end{align*}
Two further profiles measure simultaneous changes of distinct
predicates:
\begin{align*}
 J_0&=\max_{\substack{q<q'\\P_q(x)=P_{q'}(x)=0}}
       |\{a:P_q(x^a)=P_{q'}(x^a)=1\}|,\\
 J_1&=\max_{\substack{q<q'\\P_q(x)=P_{q'}(x)=1}}
       |\{a:P_q(x^a)=P_{q'}(x^a)=0\}|.
\end{align*}
All maxima range over the relevant inputs as well as the displayed
indices.  An empty maximum is zero.  Thus $J_0=J_1=0$ if $H=1$.

Unlike the source's scalar sensitivity envelopes, we retain each
$S_b(q)$ individually.  This is useful because the target conditions at
every stage use the strongest child predicate, whereas different parent
predicates use different gate thresholds.  The two joint profiles remain
scalar maxima; no array indexed by pairs is needed for our certificate.

\section{Labelled regular tournaments}
\label{sec:labels}

A tournament is an orientation of every pair of distinct vertices, with
exactly one of $i\to j$ and $j\to i$.  For an integer $h\ge1$, there
is a tournament on $k=2h+1$ vertices having outdegree exactly $h$ at every
vertex: place the vertices cyclically and direct an edge from each vertex
to the next $h$ vertices.

Each directed edge $i\to j$ will receive a label $a(i,j)\in[4]$.
The label chooses one of four child inputs in row $j$.  The following
condition limits how many almost-satisfied target rows can coexist.

\begin{definition}
An edge labelling has the $\tau$-exclusion property if there is no
$I\subseteq[k]$ of cardinality $\tau$ for which labels
$m_j\in[4]$, $j\in I$, satisfy
\[
                   a(i,j)=m_j
          \quad\text{whenever }i,j\in I\text{ and }i\to j.
\]
\end{definition}

\begin{lemma}[A uniform label choice]
\label{lem:labels}
Let $e\ge1$, $h=2^e$, $k=2h+1$, and $\tau=e+6$.  Every tournament
on $k$ vertices has a labelling by $[4]$ with the $\tau$-exclusion
property.
\end{lemma}

\begin{proof}
Independently label each directed edge uniformly from $[4]$.  For fixed
$I$ and fixed choices $(m_j)_{j\in I}$, the probability of all required
equalities is $4^{-\binom{\tau}{2}}$.  There are at most
$k^\tau4^\tau$ choices to consider.  Since $k<2^{e+2}$, the probability
of any violation is at most
\[
 k^\tau4^{\tau-\binom{\tau}{2}}
 <2^{\tau(e+2)+2\tau-\tau(\tau-1)}
 =2^{\tau(e+5-\tau)}=2^{-\tau}<1.
\]
If $\tau>k$, there are no such subsets and the conclusion is automatic.
In all cases an admissible labelling exists.
\end{proof}

This proof also defines a finite deterministic search: enumerate all
labellings and test the exclusion property until an admissible one is
found.  The probability argument proves termination.  The exponent
certificate is independent of which admissible labelling is selected.
Nothing in the proof requires enumerating the labellings during the
integer calculation supplied with this article.

\section{The nested predicates}
\label{sec:construction}

Fix a depth $d\ge1$, a positive number $b_0$ of initial blocks, and
integers $e_1,\ldots,e_d\ge1$.  Set
\[
 L=d+2,\quad H_\ell=d+1-\ell,\quad
 h_\ell=2^{e_\ell},\quad k_\ell=2h_\ell+1,\quad
 \tau_\ell=e_\ell+6.
\]
At each level choose a regular tournament on $[k_\ell]$ and an
admissible labelling from Lemma~\ref{lem:labels}.  These choices are
fixed thereafter.

\subsection{Initial predicates}

Partition $[Lb_0]$ into $b_0$ disjoint blocks $E_1,\ldots,E_{b_0}$,
each of cardinality $L$.  Let $z_j$ be the indicator of $E_j$ and put
\[
 D(x)=\min_{j\in[b_0]}\dist(x,z_j),\qquad
 P_{0,q}(x)=\one\{D(x)\le H_0-q\},\quad 1\le q\le H_0.
\]
Here $\dist$ denotes Hamming distance.  The radii range from $d$ to
zero, and distinct centres have distance $2L=2d+4$.

\begin{lemma}[Initial profiles]
\label{lem:base}
The initial predicates are nested, vanish at zero, and satisfy
\begin{gather}
 S_{0,0}(q)\le L-q,\qquad
 S_{0,1}(q)\le Lb_0-d-1+q,
 \label{eq:base}\\
 J_{0,0}=J_{0,1}=0.
 \label{eq:base-joint}
\end{gather}
Flipping any one of $E_1,\ldots,E_{b_0}$ at zero makes every
initial predicate equal to one.
\end{lemma}

\begin{proof}
Nesting follows by decreasing the radius.  At zero, $D(\zero)=L>d$,
while $D(z_j)=0$, proving the claims about zero and the blocks.

Fix a radius $R=H_0-q$.  At a rejecting input, a successful one-bit
flip must move toward a centre at distance exactly $R+1$.  There is at
most one centre within distance $d+1$ of that input: two would have
mutual distance at most $2d+2$, contrary to their separation.  There
are at most $R+1=L-q$ possible successful flips.  At an accepting
input lying strictly inside an $R$-ball, no bit flip can leave that
ball.  On its boundary, at most $Lb_0-R=Lb_0-d-1+q$ flips increase
the distance from its centre.  These observations prove
Equation~\eqref{eq:base}.

The function $D$ is $1$-Lipschitz on edges of the Boolean cube, by the
triangle inequality.  To cross two distinct integer-radius thresholds
in the same direction would require changing $D$ by at least two.
This proves Equation~\eqref{eq:base-joint}.
\end{proof}

\subsection{One recursive stage}

At level $\ell\ge1$, split the input into $4k_\ell$ disjoint child
inputs, arranged in rows $j\in[k_\ell]$ and positions $c\in[4]$.
Each child is an input to the preceding level.  Write it as $x_{j,c}$,
and abbreviate
\[
                 H=H_\ell,\qquad Q=H+1=H_{\ell-1}.
\]
For each $q\in[H]$, clause $i$ requires
\begin{align}
 P_{\ell-1,Q}(x_{i,c})&=1
                 &&\text{for every }c\in[4],
 \label{eq:targets}\\
 P_{\ell-1,Q-q}(x_{j,a(i,j)})&=0
                 &&\text{for every edge }i\to j.
 \label{eq:gates}
\end{align}
The four conditions in Equation~\eqref{eq:targets} are its targets;
the $h_\ell$ conditions in Equation~\eqref{eq:gates} are its gates.
Set $P_{\ell,q}=1$ if some clause holds, and set it to zero otherwise.

Every condition in a fixed clause uses a different child.  The targets
lie in its own row, each gate lies in a different outgoing neighbour
row, and there is only one edge from $i$ to a given $j$.  Consequently,
a one-bit flip can change at most one condition of any fixed clause.
It may change conditions of several different clauses; our upper bounds
will explicitly allow that through union bounds.

Increasing $q$ decreases the child gate index $Q-q$, so nesting of the
child predicates makes the gate condition more restrictive.  The
targets are unchanged.  The parent predicates are therefore nested as
well.  Their common input dimension is
\begin{equation}
 N_\ell=Lb_0\,4^\ell\prod_{j=1}^\ell k_j.
 \label{eq:recursive-dimension}
\end{equation}

\section{The indexed sensitivity recurrences}
\label{sec:recurrences}

In this section unprimed profiles refer to the child family with $Q$
predicates, and primed profiles to the parent family with $H=Q-1$
predicates.  Write $h$ for its tournament outdegree and $\tau$ for its
exclusion threshold.  For a fixed parent index $q$, put $g=Q-q$.

\begin{proposition}[Four coupled profile bounds]
\label{prop:recurrences}
For every parent index $q\in[H]$,
\begin{align}
 S'_1(q)&\le hS_0(g)+4S_1(Q),
 \label{eq:one-rec}\\
 S'_0(q)&\le (\tau-1)S_0(Q)
                  +\max\{S_1(g)+J_1,\,3J_1\}.
 \label{eq:zero-rec}
\end{align}
The joint profiles satisfy
\begin{align}
 J'_1&\le hJ_0+4S_1(Q),
 \label{eq:joint-one-rec}\\
 J'_0&\le (\tau-1)S_0(Q)+3J_1.
 \label{eq:joint-zero-rec}
\end{align}
\end{proposition}

\begin{proof}
\emph{One-to-zero transitions.}
At an input where $P'_q=1$, fix a satisfied clause.  Every flip that
makes $P'_q=0$ must break that same clause.  Breaking one of its $h$
gates changes child predicate $g$ from zero to one, while breaking one
of its four targets changes child predicate $Q$ from one to zero.
Equation~\eqref{eq:one-rec} follows.

\emph{Candidate clauses at a rejecting input.}
Fix an input $x$ with $P'_q(x)=0$.  Let $\mathcal T$ be the clauses
that fail exactly one target and no gate at $x$; let $\mathcal G$ be
the clauses that fail exactly one gate and no target.  These lists are
chosen before the bit flip.  If one bit makes a clause hold, that
clause must belong to one of the two lists, because a bit can affect
at most one of its conditions.

For $j\in\mathcal T$, write $m_j$ for the position of its unique
failed target.  If $i,j\in\mathcal T$ and $i\to j$, clause $i$ has
a passing gate in child $(j,a(i,j))$.  Its predicate $g$ is zero, so
its strongest predicate $Q$ is also zero.  Row $j$ has just one such
failed target, forcing $a(i,j)=m_j$.  Thus any $\tau$ vertices of
$\mathcal T$ would violate the exclusion property, and
\[
                         |\mathcal T|\le\tau-1.
\]
Repairing a target in this list changes predicate $Q$ from zero to
one.  The total contribution is at most $(\tau-1)S_0(Q)$.

All targets in the row of a vertex of $\mathcal G$ are true.  Hence an
edge $i\to j$ inside $\mathcal G$ supplies a failed gate of clause
$i$: its selected child in row $j$ has predicate $Q=1$, and therefore
$g=1$.  Every vertex of the induced tournament on $\mathcal G$ has
outdegree at most one.  If $u=|\mathcal G|$, edge counting gives
\[
                         \frac{u(u-1)}2\le u,
                         \qquad u\le3.
\]

Suppose a vertex of $\mathcal G$ has an outgoing edge within
$\mathcal G$.  That edge is its unique failed gate.  Initially the
selected child has both predicates $g$ and $Q$ equal to one.  A repair
makes $g=0$ and hence $Q=0$.  Since $g<Q$, this contributes at most
$J_1$ bit flips.

If $u=3$, all three induced outdegrees are one, so every gate repair
has this joint bound and the total is at most $3J_1$.  If $u=2$, one
vertex is the source of the sole induced edge and the other is a sink.
The first contributes at most $J_1$; the sink's failed gate contributes
at most $S_1(g)$.  For $u=1$ the bound is $S_1(g)$, and for $u=0$
it is zero.  In all cases the gate contribution is at most
\[
                         \max\{S_1(g)+J_1,3J_1\}.
\]
Adding the target contribution proves Equation~\eqref{eq:zero-rec}.
This case split is the sharpened ordinary-sensitivity gate count.

\emph{Joint transitions from one.}
Fix $q<q'$ and an input where both parent predicates are one.  Choose
a clause that holds for $q'$.  By nesting it also holds for $q$.
Let $g=Q-q$ and $g'=Q-q'$, so $g'<g<Q$.  A flip that makes both
parent predicates zero must break this fixed clause even for $q$.
A target break contributes at most $4S_1(Q)$ in total.  Otherwise a
gate must change from passing for $q'$ to failing for $q$: it begins
with $g'=g=0$, and finishes with $g=g'=1$.  Such a gate has at most
$J_0$ repairs in this direction, and there are $h$ gates.  Taking
maxima proves Equation~\eqref{eq:joint-one-rec}.

\emph{Joint transitions from zero.}
Fix $q<q'$ and an input where both parents are zero.  Form the two
candidate lists for the weaker parent $q$.  For a flip making both
parents one, choose a clause that holds afterward for $q'$.  It also
holds afterward for $q$, so it belongs to those fixed initial lists.
Its choice is allowed to depend on the flipped bit.

If it is in $\mathcal T$, the preceding target count applies.  If it
is in $\mathcal G$, its unique failed gate at index $q$ has $g=1$
initially, which implies $g'=1$.  To hold for $q'$ after the flip,
that gate must have $g'=0$, which implies $g=0$.  Thus each such
clause contributes at most $J_1$, including a possible sink.  There
are at most three clauses in $\mathcal G$, proving
Equation~\eqref{eq:joint-zero-rec}.

The argument only counts possible successful coordinates by unions of
sets.  Overlap among the coordinates counted for different clauses can
only improve the upper bound.  When $H=1$, the joint maxima are empty
and their displayed inequalities remain valid.
\end{proof}

\subsection{Numerical comparison sequences}

The finite calculation uses upper envelopes rather than exact
sensitivities.  Initialize, for $q\in[d+1]$,
\begin{equation}
 A_{0,q}=L-q,\qquad C_{0,q}=Lb_0-d-1+q,
 \qquad U_0=V_0=0.
 \label{eq:initial-envelopes}
\end{equation}
At level $\ell$, set $H=H_\ell$, $Q=H+1$, and, for $q\in[H]$,
define
\begin{align}
 A_{\ell,q}
  &=(\tau_\ell-1)A_{\ell-1,Q}
       +\max\{C_{\ell-1,Q-q}+V_{\ell-1},\,3V_{\ell-1}\},
 \label{eq:A}\\
 C_{\ell,q}
  &=h_\ell A_{\ell-1,Q-q}+4C_{\ell-1,Q},
 \label{eq:C}\\
 U_\ell
  &=(\tau_\ell-1)A_{\ell-1,Q}+3V_{\ell-1},
 \label{eq:U}\\
 V_\ell
  &=h_\ell U_{\ell-1}+4C_{\ell-1,Q}.
 \label{eq:V}
\end{align}
Every quantity here is a nonnegative integer.  The joint envelopes may
remain positive at the final level even though the actual joint profiles
are zero; this is harmless because they are only upper bounds.

\begin{corollary}[Validity of the envelopes]
\label{cor:envelopes}
At every level and valid predicate index,
\[
 S_{\ell,0}(q)\le A_{\ell,q},\quad
 S_{\ell,1}(q)\le C_{\ell,q},\quad
 J_{\ell,0}\le U_\ell,\quad J_{\ell,1}\le V_\ell.
\]
\end{corollary}

\begin{proof}
The base case is Lemma~\ref{lem:base}.  All sums, scalar multiples,
and maxima in Equations~\eqref{eq:A}--\eqref{eq:V} are nondecreasing
functions of their arguments.  Proposition~\ref{prop:recurrences}
therefore proves the assertion inductively.
\end{proof}

The separate index information is used throughout this induction.  In
particular, replacing $A_{\ell-1,Q}$ by the largest value of
$A_{\ell-1,q}$ can lose useful information.  The finite certificate
never makes that replacement.

\section{Disjoint sensitive blocks and final balancing}
\label{sec:blocks}

\begin{lemma}[Blocks shared by all predicates]
\label{lem:blocks}
At level $\ell$, every predicate vanishes at zero.  There are
\[
             b_\ell=b_0\prod_{j=1}^\ell k_j
\]
pairwise disjoint nonempty coordinate blocks, each of cardinality
$L4^\ell$, such that flipping any one of them at zero makes every
predicate at that level equal to one.
\end{lemma}

\begin{proof}
Lemma~\ref{lem:base} supplies the base case.  At a higher level all
targets are false at zero, so all clauses fail.  Enumerate the
preceding-level block family as $B_1,\ldots,B_{b_{\ell-1}}$.
For each row $i$ and each index $j$, take the union of the copies of
$B_j$ in all four children of row $i$.  Flipping this union makes
every target in clause $i$ true.  All its gates lie in other rows and
are still zero, so that clause holds for every parent index.

Different indices $j$ give disjoint unions within a row, because the
child block families are disjoint; different rows use disjoint input
coordinates.  The new family therefore has $k_\ell b_{\ell-1}$
members, each four times the previous size.
\end{proof}

At depth $d$, only one predicate remains.  Let $f$ be the OR of $m_0$
disjoint copies of $P_{d,1}$.  This OR is the standard balancing step
used in sensitivity constructions; we spell out the estimates to avoid
any implicit equality between the sensitivities at zero and at one.

\begin{proposition}[OR balancing]
\label{prop:or}
The resulting total Boolean function is nonconstant, vanishes at zero,
and satisfies
\begin{equation}
 \sens(f)\le\max\{m_0A_{d,1},C_{d,1}\},\qquad
 \bs(f,\zero)\ge m_0b_0\prod_{j=1}^d k_j.
 \label{eq:OR}
\end{equation}
Its input dimension is $m_0N_d$.
\end{proposition}

\begin{proof}
At a rejecting OR input, every copy rejects; possible successful bit
flips lie in their disjoint inputs and contribute at most
$m_0A_{d,1}$.  At an accepting OR input with at least two accepting
copies, no one-bit flip can make the OR reject.  If exactly one copy
accepts, every sensitive bit belongs to that copy, contributing at most
$C_{d,1}$.  These are upper bounds and do not assume that either
maximum is attained simultaneously with another.

Each copy contributes the disjoint blocks from Lemma~\ref{lem:blocks}
at zero; the families in distinct copies are also disjoint.  Flipping
one of these blocks makes the OR accept.  This gives the block lower
bound, and also proves nonconstancy.
\end{proof}

\section{A twenty-stage integer certificate}
\label{sec:certificate}

We now fix
\begin{equation}
 d=20,\qquad L=22,\qquad b_0=2^{18},\qquad m_0=100.
 \label{eq:parameters}
\end{equation}
Every stage has four children per row, $h_\ell=2^{e_\ell}$,
$k_\ell=2^{e_\ell+1}+1$, and $\tau_\ell=e_\ell+6$.
The exponent schedule is given in Table~\ref{tab:parameters}.
Its entries sum to $608$.

\begin{table}[htbp]
\centering
\begin{tabular}{@{}rrrrrr@{}}
\toprule
Stage $\ell$ & $e_\ell$ & $\tau_\ell$
  & Stage $\ell$ & $e_\ell$ & $\tau_\ell$\\
\midrule
1 &43&49&11&30&36\\
2 &43&49&12&28&34\\
3 &39&45&13&26&32\\
4 &39&45&14&24&30\\
5 &40&46&15&23&29\\
6 &40&46&16&22&28\\
7 &37&43&17&20&26\\
8 &36&42&18&19&25\\
9 &34&40&19&18&24\\
10&32&38&20&15&21\\
\bottomrule
\end{tabular}
\caption{The fixed schedule.  All label-existence conditions follow from
Lemma~\ref{lem:labels}; no numerical probability estimate is needed.}
\label{tab:parameters}
\end{table}

\begin{lemma}[Finite arithmetic certificate]
\label{lem:arithmetic}
For Equations~\eqref{eq:initial-envelopes}--\eqref{eq:V} and the
parameters of Table~\ref{tab:parameters},
\begin{equation}
       100A_{20,1}<9\cdot2^{318},\qquad
       C_{20,1}<9\cdot2^{318}.
 \label{eq:finite-inequalities}
\end{equation}
\end{lemma}

\begin{proof}[Proof by exact finite computation]
Appendix~\ref{app:verifier} contains the complete integer computation.
Its arrays $a$ and $c$ store $A_{\ell,q}$ and $C_{\ell,q}$ in
increasing order of $q$, with zero-based array indexing.  Its variables
$u,v$ store $U_\ell,V_\ell$.  At each stage, $H$ is one less than the
incoming array length, so the strongest child entry is the last one;
the child index $Q-q$ is stored at array index $H-q$.  The program is
therefore a direct evaluation of the four displayed recurrences,
not an approximation to them.

The initialization stores Equation~\eqref{eq:initial-envelopes}.
There are exactly twenty stages and every operation is addition,
multiplication, an integer power, comparison, or maximum on integers.
The two concluding integer comparisons are precisely
Equation~\eqref{eq:finite-inequalities}.  The accompanying executable
\texttt{verify.py} performs the same calculation, checks additional
dimension and label identities, and prints the final exact values.
No floating-point number, random sample, external solver, or third-party
package is used in the certificate.
\end{proof}

\begin{remark}[Role of the computation]
The search that suggested the schedule is not part of the proof.  Any
method of finding the twenty exponents would give the same theorem once
the fixed schedule and its inequalities were certified.  Conversely,
checking the integer computation alone would not prove the theorem
without Proposition~\ref{prop:recurrences} and the block construction.
The logical division is explicit: those are general combinatorial
proofs; Lemma~\ref{lem:arithmetic} is a finite arithmetic calculation.
\end{remark}

By Proposition~\ref{prop:or} and Lemma~\ref{lem:arithmetic}, the seed
satisfies $\sens(f)\le S_*$.  Meanwhile,
\begin{align*}
 \bs(f,\zero)
 &\ge 100\cdot2^{18}
            \prod_{\ell=1}^{20}(2^{e_\ell+1}+1)\\
 &>100\cdot2^{18+608+20}
   =100\cdot2^{646}=B_*.
\end{align*}
Its specified input dimension is
\begin{equation}
 n_*=100\cdot22\cdot2^{18}\cdot4^{20}
             \prod_{\ell=1}^{20}(2^{e_\ell+1}+1).
 \label{eq:seed-dimension}
\end{equation}
The disjoint sensitive blocks occupy precisely these inputs when one
counts their common size $22\cdot4^{20}$.  This provides a useful
consistency check on the construction's bookkeeping.

The rational exponent comparison has a short elementary certificate.
Directly,
\[
 \frac{B_*^{32}}{S_*^{65}}
 =\frac{4\cdot100^{32}}{9^{65}}>1.
\]
Indeed $100>81$ and
$2\cdot100^3=2{,}000{,}000>1{,}594{,}323=3\cdot81^3$.
Multiplication by $(100/81)^{13}>1$ gives
$2\cdot100^{16}>3\cdot81^{16}=3^{65}$; squaring proves
$4\cdot100^{32}>9^{65}$.  Since $S_*>1$, this establishes
\begin{equation}
                    \frac{\log B_*}{\log S_*}>\frac{65}{32}.
 \label{eq:exponent-exact}
\end{equation}
The decimal expansion in the abstract is informational.  It is not
used in Equation~\eqref{eq:exponent-exact} or in any theorem.

\section{Composition and the infinite family}
\label{sec:composition}

The composition principle is standard \cite{Tal2013}; for completeness,
we give precisely the forms required at zero.

\begin{lemma}[Disjoint composition]
\label{lem:composition}
Let $f:\bits^a\to\bits$ and $g:\bits^b\to\bits$ be total Boolean
functions, and define
\[
 (f\circ g)(x_1,\ldots,x_a)
          =f(g(x_1),\ldots,g(x_a)),\qquad x_i\in\bits^b.
\]
Then
\[
                     \sens(f\circ g)\le\sens(f)\sens(g).
\]
If $g(\zero)=0$, then $(f\circ g)(\zero)=f(\zero)$ and
\[
                  \bs(f\circ g,\zero)
                  \ge\bs(f,\zero)\bs(g,\zero).
\]
\end{lemma}

\begin{proof}
At an input $x=(x_1,\ldots,x_a)$ let
$y=(g(x_1),\ldots,g(x_a))$.  A raw bit in child $x_i$ is sensitive
for the composition exactly when that bit is sensitive for $g$ at
$x_i$ and coordinate $i$ is sensitive for $f$ at $y$.  Thus
\[
 \sens(f\circ g,x)
    =\sum_{i:f(y^i)\ne f(y)}\sens(g,x_i)
    \le\sens(f,y)\sens(g)
    \le\sens(f)\sens(g).
\]

For the second assertion, let $C_1,\ldots,C_p$ and
$D_1,\ldots,D_q$ be disjoint sensitive blocks at zero for $f$ and
$g$.  Index the composed input coordinates by $[a]\times[b]$.
The sets $C_i\times D_j$ are nonempty and pairwise disjoint.  Flipping
one of them at zero makes exactly the children indexed by $C_i$
output one, because $g(\zero)=0$ and $D_j$ is sensitive there.  The
outer input becomes $\zero^{C_i}$, which changes $f$.  There are
$pq$ such blocks, proving the claim.
\end{proof}

\begin{proof}[Proof of Theorem~\ref{thm:main}]
Take the seed $f$ of Sections~\ref{sec:blocks}--\ref{sec:certificate}.
Set $F_1=f$ and $F_{m+1}=f\circ F_m$, using disjoint copies of $F_m$.
These functions have $n_*^m$ inputs and vanish at zero.
Lemma~\ref{lem:composition} gives
\[
                  \sens(F_m)\le S_*^m,\qquad
                  \bs(F_m,\zero)\ge B_*^m.
\]
The block bound proves nonconstancy and tends to infinity.  A
nonconstant Boolean function has positive sensitivity: along a path
of one-bit flips between two inputs with different outputs, at least
one edge changes the output.

Since $S_*^{\alpha_*}=B_*$,
\[
 \sens(F_m)^{\alpha_*}
 \le S_*^{m\alpha_*}=B_*^m\le\bs(F_m,\zero).
\]
Equation~\eqref{eq:exponent-exact} gives $\alpha_*>65/32$.  More
explicitly,
\[
 \frac{\bs(F_m,\zero)}{\sens(F_m)^{65/32}}
 \ge\left(\frac{B_*}{S_*^{65/32}}\right)^m,
\]
and the base on the right is strictly greater than one.  This proves
the asserted divergence and completes the theorem.
\end{proof}

\section{What the refinement changes and what remains open}
\label{sec:research}

The source architecture and its central insight are retained: nesting
forces some otherwise expensive gate repairs to change two child
predicates simultaneously.  The initial joint profiles vanish, and a
controlled number of recursive stages exploits that asymmetry before
composition amplifies the resulting seed.  The refinement has three
specific ingredients.

First, the outdegrees $h_\ell$ are allowed to vary by level.  The
schedule begins with very large tournaments and ends with much smaller
ones.  Each stage still uses the same proved local construction; the
source's global requirement that one size parameter serve every level
is unnecessary for the finite certificate.

Second, each sensitivity $S_b(q)$ is tracked separately.  The strongest
target index is used exactly where the proof requires it, and the gate
index remains $Q-q$.  This preserves information lost by taking an
early maximum over all indices.  The joint profiles can remain scalar,
so this gain does not require an exponentially large collection of
auxiliary quantities.

Third, the gate candidate list is counted according to its size.  A
three-vertex list is a directed cycle and has no sink; a two-vertex
list has one sink.  The resulting maximum
$\max\{S_1(g)+J_1,3J_1\}$ improves the simultaneous sum
$S_1(g)+3J_1$ used in the source's convenient upper bound.  All three
changes are inequalities for every relevant input, not assumptions
about the inputs where the displayed sensitive blocks occur.

The following questions are suitable next steps.

\begin{enumerate}
\item \textbf{Optimize the certified exponent.}
For fixed depth and a prescribed input-size budget, optimize the integer
schedule $(e_\ell)$, the number of initial blocks, and the final OR
multiplicity under Equations~\eqref{eq:A}--\eqref{eq:V}.  More
generally, vary the number of positions per row and use its corresponding
label-exclusion condition.  A search result should always be converted
to a finite exact certificate, as here.

\item \textbf{Determine the accuracy of the envelopes.}
For small parameters, compare actual profile maxima with the upper
bounds.  Which case of the target/gate analysis is attainable
simultaneously?  Strong incompatibilities between the candidate target
and gate lists could improve the recurrence rather than merely its
parameters.  Small exhaustive calculations are evidence for such a
lemma, not a substitute for its proof at the large certified parameters.

\item \textbf{Use pair-specific joint profiles.}
For parent indices $q<q'$, a nonsink gate repair at zero changes both
the child predicate $Q-q'$ and the strongest child predicate $Q$.
Only a possible sink needs the pair $Q-q',Q-q$.  Tracking those pairs
separately gives additional valid inequalities.  Whether they yield a
material exponent improvement after optimization is open in this note;
the present certificate does not need them.

\item \textbf{Reduce the label-search and seed size.}
The union bound proves existence but is inefficient.  Explicit
labellings, conditional-expectation constructions, or a local-lemma
analysis may permit smaller tournaments or exclusion thresholds.
The meaningful target is a reduction in the seed input dimension while
preserving a proved separation exponent.

\item \textbf{Identify an architectural improvement.}
The factor three in the joint gate bound comes from a possible directed
three-cycle among almost-satisfied clauses.  A different clause system
that retains a large supply of disjoint sensitive blocks while
restricting this obstruction may improve the exponent more than further
numerical optimization.

\item \textbf{Formalize the generalized recurrence.}
Extend the repository's sensitivity definitions and proof development to
indexed sensitivities and level-dependent tournament parameters.  Once
the general combinatorial lemmas are formalized, the fixed numerical
certificate should become a small theorem about integer recurrences.
The comparator statement should use the same ordinary-sensitivity and
block-sensitivity definitions as the prose theorem.
\end{enumerate}

These questions distinguish optimization of a sufficient certificate
from improvements to the Boolean functions themselves.  The latter may
require new information about which worst-case profiles can coexist.
Neither the finite search that found our schedule nor the upper
envelopes prove a barrier for the underlying construction.

\appendix
\section{Complete core of the integer verifier}
\label{app:verifier}

The following program is sufficient for the finite inequalities.  The
packaged \texttt{verify.py} adds explanatory output and consistency
checks but uses these same recurrences.  Python's integers have arbitrary
precision; their use here involves no rounding or probabilistic step.
The explicit checks remain active when Python is run with optimization
enabled.

\begingroup
\small
\begin{verbatim}
def require(condition, message):
    if not condition:
        raise ArithmeticError(message)

E = [43,43,39,39,40,40,37,36,34,32,
     30,28,26,24,23,22,20,19,18,15]
d, b0, copies = 20, 2**18, 100
L = d + 2
a = [L-q for q in range(1, d+2)]
c = [L*b0-d-1+q for q in range(1, d+2)]
u = v = 0

for e in E:
    h, tau = 2**e, e+6
    H = len(a)-1
    new_a = [
        (tau-1)*a[-1] + max(c[H-q]+v, 3*v)
        for q in range(1, H+1)
    ]
    new_c = [
        h*a[H-q] + 4*c[-1]
        for q in range(1, H+1)
    ]
    new_u = (tau-1)*a[-1] + 3*v
    new_v = h*u + 4*c[-1]
    a, c, u, v = new_a, new_c, new_u, new_v

require(len(a) == len(c) == 1, "Final profile lengths")
require(sum(E) == 608, "Sum of outdegree exponents")
S_star = 9 * 2**318
B_star = 100 * 2**646
require(copies*a[0] < S_star, "Zero-sensitivity bound")
require(c[0] < S_star, "One-sensitivity bound")
require(B_star**32 > S_star**65, "Exponent exceeds 65/32")
print("All exact certificate inequalities hold.")
\end{verbatim}
\endgroup

\begin{thebibliography}{9}

\bibitem{OpenAI2026}
\begingroup
\raggedright
OpenAI.
\emph{A superquadratic separation between sensitivity and block
sensitivity}. Repository manuscript, September 25, 2026.
Reviewed in \texttt{VladimirReshetnikov/math}, commit
\texttt{adc7f1241b42e322a6451854ab7e4b4c146bf78a},
October 6, 2026, 21:58:50 UTC.
\href{https://github.com/VladimirReshetnikov/math/tree/adc7f1241b42e322a6451854ab7e4b4c146bf78a/preprints/A-superquadratic-separation-between-sensitivity-and-block-sensitivity-September-25-2026/build}{Pinned manuscript source}.
Source sections: \texttt{00-introduction.tex},
\texttt{01-construction.tex}, \texttt{02-recurrences.tex}, and
\texttt{03-separation.tex}, under \texttt{build/sections/}.
The related proof entry is
\href{https://github.com/VladimirReshetnikov/math/blob/adc7f1241b42e322a6451854ab7e4b4c146bf78a/lean/OAI/Combinatorics/Sensitivity/Separation.lean}{\texttt{OAI.Combinatorics.Sensitivity.Separation}}.
\par
\endgroup

\bibitem{Huang2019}
Hao Huang.
Induced subgraphs of hypercubes and a proof of the Sensitivity Conjecture.
\emph{Annals of Mathematics} \textbf{190} (2019), 949--955.
\href{https://doi.org/10.4007/annals.2019.190.3.6}{doi:10.4007/annals.2019.190.3.6}.

\bibitem{Meiburg2026}
Alexander Meiburg.
\emph{Block Sensitivity can exceed Spectral Sensitivity Squared}.
Preprint, August 1, 2026.
\href{https://arxiv.org/abs/2608.00851}{arXiv:2608.00851}.
The spectral measure studied there is distinct from the ordinary
sensitivity used in this companion.

\bibitem{Tal2013}
Avishay Tal.
Properties and Applications of Boolean Function Composition.
In \emph{Proceedings of ITCS 2013}, 441--454.
\href{https://doi.org/10.1145/2422436.2422485}{doi:10.1145/2422436.2422485}.

\end{thebibliography}
\end{document}
